\clearpage
\SourceLeaf{495}%
\fancyhead[L]{\small INHALT.}
\pdfbookmark[0]{Inhalt.}{inhalt}
\begin{center}\Large INHALT.\par\bigskip GAUSS WERKE BAND IV.\par\medskip\large WAHRSCHEINLICHKEITSRECHNUNG UND GEOMETRIE.\end{center}
\begingroup\fontsize{10.5}{13}\selectfont\interlinepenalty=10000
\par\medskip\Needspace{5\baselineskip}
\begin{center}WAHRSCHEINLICHKEITSRECHNUNG.\end{center}
\par\medskip\Needspace{5\baselineskip}
\noindent\textit{Abhandlungen.}\par
\par\noindent\hangindent=1em Theoria combinationis observationum erroribus minimis obnoxiae. Pars prior \hfill\penalty0\hbox{}\nobreak\hfill\mbox{1821 Febr. 15 Seite 1}\par
\par\noindent\hangindent=1em Pars posterior . . . . . . . . . . . . . . . . . . . . .\hfill\penalty0\hbox{}\nobreak\hfill\mbox{1823 Febr. 2 — 27}\par
\par\noindent\hangindent=1em Supplementum theoriae combinationis observationum erroribus minimis obnoxiae . . . . . . . . . . . . . . . . . . . . . .\hfill\penalty0\hbox{}\nobreak\hfill\mbox{1826 Sept. 16 — 55}\par
\par\medskip\Needspace{5\baselineskip}
\noindent\textit{Anzeigen eigner Abhandlungen.}\par
\par\noindent\hangindent=1em Theoria combinationis observationum erroribus minimis obnoxiae. Pars prior \hfill\penalty0\hbox{}\nobreak\hfill\mbox{1821 Febr. 26 — 95}\par
\par\noindent\hangindent=1em Pars posterior . . . . . . . . . . . . . . . . . . . . .\hfill\penalty0\hbox{}\nobreak\hfill\mbox{1823 Febr. 24 — 100}\par
\par\noindent\hangindent=1em Supplementum theoriae combinationis observationum etc. . . . . . .\hfill\penalty0\hbox{}\nobreak\hfill\mbox{1826 Sept. 25 — 104}\par
\par\medskip\Needspace{5\baselineskip}
\noindent\textit{Aufsatz.}\par
\par\noindent\hangindent=1em Bestimmung der Genauigkeit der Beobachtungen . . . . . . . . .\hfill\penalty0\hbox{}\nobreak\hfill\mbox{1816 März . . — 109}\par
\par\medskip\Needspace{5\baselineskip}
\noindent\textit{Nachlass.}\par
\par\noindent\hangindent=1em Anwendung der Wahrscheinlichkeitsrechnung auf die Bestimmung der Bilanz für Witwenkassen . . . . . . . . . . . . . . . . . .\hfill\penalty0\hbox{}\nobreak\hfill\mbox{1845 Jan. 9 — 119}\par
\par\noindent\hangindent=1em Tafeln zur Bestimmung von Leibrenten . . . . . . . . . . . .\hfill\penalty0\hbox{}\nobreak\hfill\mbox{— 174}\par
\par\noindent\hangindent=1em Einrichtung und Gebrauch dieser Tafeln . . . . . . . . . . . .\hfill\penalty0\hbox{}\nobreak\hfill\mbox{— 184}\par
\par\medskip\Needspace{5\baselineskip}
\begin{center}GEOMETRIE.\end{center}
\par\medskip\Needspace{5\baselineskip}
\noindent\textit{Abhandlungen.}\par
\par\noindent\hangindent=1em Allgemeine Auflösung der Aufgabe die Theile einer gegebenen Fläche auf einer andern gegebenen Fläche so abzubilden, dass die Abbildung dem Abgebildeten in den kleinsten Theilen ähnlich wird . . . . . . .\hfill\penalty0\hbox{}\nobreak\hfill\mbox{1822 Dec. . . — 189}\par
\par\noindent\hangindent=1em Disquisitiones generales circa superficies curvas . . . . . . . . . .\hfill\penalty0\hbox{}\nobreak\hfill\mbox{1827 Oct. 8 — 217}\par
\par\noindent\hangindent=1em Untersuchungen über Gegenstände der höhern Geodaesie. Erste Abhandlung \hfill\penalty0\hbox{}\nobreak\hfill\mbox{1843 Oct. 23 — 259}\par
\par\noindent\hangindent=1em Zweite Abhandlung . . . . . . . . . . . . . . . . . . .\hfill\penalty0\hbox{}\nobreak\hfill\mbox{1846 Sept. 1 — 301}\par
\endgroup
